% Иллюстрация к выводу формулы Стирлинга:
% сумма ln k · 1 как сумма площадей прямоугольников под графиком y = ln x,
% приближённо равная интегралу от 1 до n.
% Компиляция: pdflatex stirling-riemann-figure.tex
\documentclass[tikz,border=8pt]{standalone}
\usepackage[T2A]{fontenc}
\usepackage[utf8]{inputenc}
\usepackage[russian]{babel}
\usetikzlibrary{arrows.meta}

\begin{document}
\begin{tikzpicture}[
    >={Stealth[length=2.4mm]},
    axis/.style  = {->, semithick},
    curve/.style = {very thick, blue!55!black},
    rect/.style  = {semithick, red!70!black, fill=red!12},
  ]
  % масштаб по y: 1.6*ln(x)

  % ================= прямоугольники: [k, k+1] высоты ln k =================
  % (прямоугольник для k=1 имеет высоту ln 1 = 0 и не виден)
  \draw[rect] (2,0) rectangle (3,1.109);
  \draw[rect] (3,0) rectangle (4,1.758);
  \draw[rect] (4,0) rectangle (5,2.218);
  \draw[rect] (5,0) rectangle (6,2.575);
  \draw[rect] (6,0) rectangle (7,2.867);
  \draw[rect] (7,0) rectangle (8,3.113);
  \node[red!70!black] at (4.5,1.15) {$\ln k$};

  % ================= оси =================
  \draw[axis] (-0.5,0) -- (9.4,0) node[below right=-2pt] {$x$};
  \draw[axis] (0,-0.6) -- (0,3.9) node[above] {$y$};

  % отметки на оси x
  \draw (1,-0.07) -- (1,0.07);
  \node[below] at (1.07,-0.19) {$1$};
  \foreach \x/\lbl in {2/2, 3/3, 4/k, 5/k+1, 7/n-1, 8/n} {
    \draw (\x,-0.07) -- (\x,0.07);
    \node[below] at (\x,-0.19) {$\lbl$};
  }
  \draw (6,-0.07) -- (6,0.07);
  \node[below] at (6,-0.19) {$\dots$};

  % ================= график логарифма =================
  \draw[curve, domain=0.85:8.7, smooth, samples=120]
        plot (\x, {1.6*ln(\x)});
  \node[blue!55!black] at (7.3,3.62) {$y=\ln x$};

  % точки касания: верхний левый угол прямоугольника лежит на графике
  \foreach \x in {1,2,...,8} {
    \fill[blue!55!black] (\x,{1.6*ln(\x)}) circle (0.05);
  }

\end{tikzpicture}
\end{document}
